% ats-single-column — matching cover letter.
% The renderer substitutes %%PDFMETA%%, %%HEADER%% and %%BODY%%.

\documentclass[11pt,a4paper]{article}

\usepackage[a4paper,top=1.8cm,bottom=1.8cm,left=1.9cm,right=1.9cm]{geometry}
\usepackage{fontspec}
\setmainfont{texgyreheros}[
  Extension      = .otf,
  UprightFont    = *-regular,
  BoldFont       = *-bold,
  ItalicFont     = *-italic,
  BoldItalicFont = *-bolditalic,
]
\usepackage{xcolor}
\usepackage{microtype}
\usepackage[hidelinks]{hyperref}

\definecolor{rulegray}{gray}{0.60}
\definecolor{muted}{gray}{0.30}

\pagestyle{empty}
\setlength{\parindent}{0pt}
\setlength{\parskip}{0.75em}
\linespread{1.08}

\newcommand{\clname}[1]{{\fontsize{17}{19}\selectfont\bfseries #1}\par\vspace{0.2em}}
\newcommand{\clcontact}[1]{{\small\color{muted}#1}\par\vspace{0.3em}%
  {\color{rulegray}\rule{\linewidth}{0.5pt}}\par}
\newcommand{\cldate}[1]{{\small\color{muted}#1}\par}

\hypersetup{
  pdftitle={Nguyen Thi Dieu Hang - Cover letter},
  pdfauthor={Nguyen Thi Dieu Hang},
  pdfcreator={},
}

\begin{document}
\clname{Nguyen Thi Dieu Hang}
\clcontact{\href{mailto:ntdieuhang192@gmail.com}{ntdieuhang192@gmail.com} $\cdot$ +84 327 840 518 $\cdot$ Ho Chi Minh City, Vietnam $\cdot$ \href{https://www.linkedin.com/in/hangnguyen118}{linkedin.com/in/hangnguyen118} $\cdot$ \href{https://github.com/hangnguyen118}{github.com/hangnguyen118}}
\cldate{05 September 2026}

Dear Hiring Team,

I am applying for the Junior Frontend Engineer role at Nexlab Technology. The posting describes Spec-Driven Development as the methodology the whole pipeline runs on, with the specification as the single source of truth and humans verifying at fixed gates. That is the process I used to build my portfolio site: I wrote the specification first, had Claude Code build to it, and validated the result automatically rather than by eye. It is a way of working I have chosen deliberately, not something I would be learning from scratch here.

My current role at VILIHA PTE. LTD. is production frontend work in Next.js — App Router and Server Components, TanStack Query and TanStack Table for client-side data, Tailwind CSS and shadcn/ui — built from detailed Figma designs in TypeScript. Before that, at Chip Nova, I clarified requirements, designed the user experience, and implemented the frontend for web and mobile, so I am used to owning a feature from specification through to shipped interface.

The part of the role I find most interesting is verification. At Chip Nova I wrote end-to-end Playwright tests to protect product quality, and I built a Playwright framework in TypeScript on the Page Object Model covering 21 test cases across four browser and device targets. That background is directly about the problem the posting names: an interface that looks right at a glance can be wrong in the state you did not open, and finding that requires deliberate checking rather than a quick look.

I should be straightforward about the gaps. I have not yet measured or optimised Core Web Vitals on a real project, set up a CI/CD pipeline, or worked with a Headless CMS or GraphQL, and my accessibility and SEO knowledge is not yet at a working level. I have about one year of professional frontend experience. I would welcome the chance to discuss how I work and to hear more about the projects the team is taking on.

Sincerely, Nguyen Thi Dieu Hang
\end{document}
